% CCSS 7.RP.A.1 — Unit Rates with Fractions
\section{Unit Rates with Fractions}

\begin{learningGoals}
\begin{itemize}
    \item Compute unit rates when quantities are fractions
    \item Compare rates by finding unit rates
\end{itemize}
\end{learningGoals}

\begin{conceptBox}{Unit Rates with Fractions}
A \textbf{unit rate} tells you how much of one quantity per \textbf{one unit} of another.
When the quantities are fractions, divide just like before:
$$\text{unit rate} = \frac{\text{first quantity}}{\text{second quantity}}$$
To divide fractions, \textbf{multiply by the reciprocal} (Keep--Change--Flip):
$$\frac{a}{b} \div \frac{c}{d} = \frac{a}{b} \times \frac{d}{c}$$
\end{conceptBox}

\begin{workedExample}{Finding a Unit Rate}
A recipe uses $\frac{3}{4}$ cup of sugar for $\frac{1}{2}$ batch. How much sugar per batch?
\solutionLabel
$\dfrac{3}{4} \div \dfrac{1}{2} = \dfrac{3}{4} \times \dfrac{2}{1} = \dfrac{6}{4} = \dfrac{3}{2}$
\answerHighlight{$1\dfrac{1}{2}$ cups per batch}
\end{workedExample}

\begin{sideBySideExample}{Comparing Two Rates}
Who walks faster?
\begin{itemize}[leftmargin=*]
    \item Ana: $\frac{2}{3}$ mi in $\frac{1}{4}$ hr
    \item Ben: $\frac{3}{4}$ mi in $\frac{1}{3}$ hr
\end{itemize}
\tcblower
\textbf{Ana:} $\frac{2}{3} \div \frac{1}{4} = \frac{8}{3} = 2\frac{2}{3}$ mph

\textbf{Ben:} $\frac{3}{4} \div \frac{1}{3} = \frac{9}{4} = 2\frac{1}{4}$ mph

$2\frac{2}{3} > 2\frac{1}{4}$ → \textbf{Ana is faster.}
\end{sideBySideExample}

\mascotSays{``Per'' means ``for every one.'' To find a unit rate, always divide!}

\begin{practiceBox}{Unit Rates with Fractions}
\resetProblems

\prob A painter covers $\frac{2}{3}$ of a wall in $\frac{1}{3}$ hour. How much wall per hour?
\answerExplain{$2$ walls per hour}{Divide the wall covered by the time: $\frac{2}{3} \div \frac{1}{3}$. Flip the second fraction and multiply: $\frac{2}{3} \times \frac{3}{1} = \frac{6}{3} = 2$ walls per hour.}

\prob A faucet leaks $\frac{1}{2}$ cup in $\frac{3}{4}$ minute. How many cups per minute?
\answerExplain{$\frac{2}{3}$ cup per minute}{Divide cups by minutes: $\frac{1}{2} \div \frac{3}{4}$. Flip and multiply: $\frac{1}{2} \times \frac{4}{3} = \frac{4}{6} = \frac{2}{3}$ cup per minute.}

\prob A car uses $\frac{5}{6}$ gallon to travel $\frac{2}{3}$ mile. How many gallons per mile?
\answerExplain{$1\frac{1}{4}$ gallons per mile}{Divide gallons by miles: $\frac{5}{6} \div \frac{2}{3}$. Flip and multiply: $\frac{5}{6} \times \frac{3}{2} = \frac{15}{12} = \frac{5}{4} = 1\frac{1}{4}$ gallons per mile.}

\prob Trail mix A: $\frac{1}{2}$ cup nuts for $\frac{1}{4}$ cup raisins. Trail mix B: $\frac{2}{3}$ cup nuts for $\frac{1}{3}$ cup raisins. Which has more nuts per cup of raisins?
\answerExplain{Equal --- both $2$ cups per cup}{Find each unit rate. Mix A: $\frac{1}{2} \div \frac{1}{4} = \frac{1}{2} \times 4 = 2$. Mix B: $\frac{2}{3} \div \frac{1}{3} = \frac{2}{3} \times 3 = 2$. Both give $2$ cups of nuts per cup of raisins, so they are equal.}

\prob A snail crawls $\frac{5}{8}$ m in $\frac{3}{4}$ hr. A beetle crawls $\frac{7}{10}$ m in $\frac{4}{5}$ hr. Which is faster?
\answerExplain{Beetle: $\frac{7}{8}$ m/hr $>$ $\frac{5}{6}$ m/hr}{Find each unit rate. Snail: $\frac{5}{8} \div \frac{3}{4} = \frac{5}{8} \times \frac{4}{3} = \frac{20}{24} = \frac{5}{6}$ m/hr. Beetle: $\frac{7}{10} \div \frac{4}{5} = \frac{7}{10} \times \frac{5}{4} = \frac{35}{40} = \frac{7}{8}$ m/hr. Since $\frac{7}{8} > \frac{5}{6}$, the beetle is faster.}

\end{practiceBox}
